\section{Further questions}
\label{sec:further-questions}

We conclude with several questions that are not settled in this paper.
Whenever $H$ is a numerical semigroup considered below, $\ell$ denotes the
least weighted degree of a nonzero homogeneous element of the defining ideal
$\fkp_H$ of $k[H]$.

\begin{question}
\label{que:seven-from-semigroups}
	Are all the seven-generated ideals $I_{\ell+1}$ arising
	from non-symmetric numerical semigroups of multiplicity four and embedding
	dimension three normal?  Is the same true for those arising from symmetric
	numerical semigroups of multiplicity five and embedding dimension three?
\end{question}

For the next question, let $H$ be a numerical
semigroup of embedding dimension three and arbitrary multiplicity.  Since
$I_\ell$ is an integrally closed monomial ideal, it is normal by
\cite[Theorem~3.1]{AtakaMatsuoka2026} whenever $\mu_S(I_\ell)\leq7$.

\begin{question}
	Suppose that $\mu_S(I_\ell)\leq7$ and
	$\mu_S(I_{\ell+1})\leq7$.  Must $I_{\ell+1}$ be normal?
\end{question}

The preceding question concerns only the passage from $I_\ell$ to
$I_{\ell+1}$.  The following formulation asks the same at every step.
Let $H$ be a numerical semigroup of embedding dimension three, and retain
the notation $S$, $\varphi_H$, and
\[
 I_h=\varphi_H^{-1}\bigl(t^hk[t]\cap k[H]\bigr)
\]
of Section~\ref{sec:target-ideals}.

\begin{question}
\label{que:successive-degree-bounds}
For $h\in H$, set
\[
 \sigma(h)=\min\{n\in H\mid n>h\}.
\]
Suppose that $\mu_S(I_h)\leq7$ and $\mu_S(I_{\sigma(h)})\leq7$.
If $I_h$ is normal, must $I_{\sigma(h)}$ be normal?
\end{question}

The following example shows that normality need not pass to the next
cutoff without restrictions on the number of generators, even when both
ideals are monomial.

\begin{example}\label{ex:successive-cutoffs-eight-generators}
Let $H=\langle6,14,21\rangle$, and give $S=k[x,y,z]$ the grading
$\deg x=6$, $\deg y=14$, and $\deg z=21$.  The defining ideal is
\[
 \fkp_H=(y^3-x^7,\ z^2-x^7),
\]
so $\ell=42$.  Since $41=6+14+21\in H$ and $42\in H$, we have
$\sigma(41)=42$.  The two successive ideals are
\[
\begin{aligned}
 I_{41}&=(x^7,x^5y,x^3y^2,y^3,x^4z,xyz,y^2z,z^2),\\
 I_{42}&=(x^7,x^5y,x^3y^2,y^3,x^4z,x^2yz,y^2z,z^2).
\end{aligned}
\]
Both are integrally closed monomial ideals with eight minimal generators.
A direct computation with the Newton polyhedron of $I_{41}^2$ shows
that $I_{41}^2$ is integrally closed.  Hence $I_{41}$ is normal by
\cite[Proposition~3.1]{ReidRobertsVitulli2003}.

On the other hand, $I_{42}=\ol{(x^7,y^3,z^2)}$ is not normal.
Indeed, put $w=x^6y^2z$.  Then $w\notin I_{42}^2$, whereas
\[
 w^2=(x^7)(x^5y)(y^3)(z^2)\in I_{42}^4.
\]
Thus $w\in\ol{I_{42}^2}\setminus I_{42}^2$.  In contrast,
$w=(xyz)(x^5y)\in I_{41}^2$.
\end{example}
